This is the core of the workshop. You will train a recommendation model
and put it into real-time service.

\vspace{0.5em}\noindent\textbf{Understand the data before you
start}\par

Amazon Personalize learns from \textbf{behaviour}, not from product
attributes. It needs a table recording who interacted with which item,
and when.

The required format has these columns, in uppercase:

\begin{verbatim}
USER_ID,ITEM_ID,EVENT_TYPE,TIMESTAMP
user-133,prod-080,view,1779328787
user-133,prod-008,view,1779328979
user-133,prod-008,add_to_cart,1779329120
user-133,prod-008,purchase,1779329340
\end{verbatim}

\texttt{TIMESTAMP} must be a Unix integer, not a date string.

\textbf{The most valuable lesson from this workshop.} Our first dataset
had 4,491 interactions generated \textbf{uniformly at random}. The
trained model performed terribly, and on reflection the reason was
obvious: collaborative filtering works by finding groups of users who
behave alike. If everyone interacts at random, no such groups exist.

Switching algorithms would not have fixed that. The data had to change.

The second dataset simulated realistic behaviour with five properties:
grouping into browsing sessions, a conversion funnel from view to
add-to-cart to purchase, a power-law distribution so a small set of
products attracts most views, time-of-day rhythms, and users split into
segments with different tastes.

Results on the same recipe:

{\def\LTcaptype{none} % do not increment counter
\begingroup
\small
\setlength{\tabcolsep}{3pt}
\renewcommand{\arraystretch}{1.15}
\begin{longtable}[]{@{}lll@{}}
\toprule\noalign{}
Metric & Random data & Behaviour-simulated data \\
\midrule\noalign{}
\endhead
\bottomrule\noalign{}
\endlastfoot
Precision@5 & 0.0889 & 0.4348 \\
NDCG@10 & 0.1799 & 0.6512 \\
MRR@25 & 0.1216 & 0.7130 \\
Coverage & 0.8218 & 0.9505 \\
\end{longtable}
\endgroup
}

\vspace{0.5em}\noindent\textbf{Step 1. Create a bucket and upload the
data}\par

\begin{Shaded}
\begin{Highlighting}[]
\ExtensionTok{aws}\NormalTok{ s3 mb s3://fcj{-}recsys{-}data{-}}\OperatorTok{\textless{}}\NormalTok{your{-}name}\OperatorTok{\textgreater{}}
\ExtensionTok{aws}\NormalTok{ s3 cp interactions.csv s3://fcj{-}recsys{-}data{-}}\OperatorTok{\textless{}}\NormalTok{your{-}name}\OperatorTok{\textgreater{}}\NormalTok{/personalize/}
\ExtensionTok{aws}\NormalTok{ s3 cp items.csv s3://fcj{-}recsys{-}data{-}}\OperatorTok{\textless{}}\NormalTok{your{-}name}\OperatorTok{\textgreater{}}\NormalTok{/personalize/}
\end{Highlighting}
\end{Shaded}

\vspace{0.5em}\noindent\textbf{Step 2. Let Personalize read the
bucket}\par

Personalize is a separate service and needs explicit permission to read
your bucket.

Create the role:

\begin{enumerate}
\def\labelenumi{\arabic{enumi}.}
\tightlist
\item
  Go to the \href{https://console.aws.amazon.com/iam/home\#/roles}{IAM
  console} and choose \textbf{Create role}
\item
  \textbf{Trusted entity type}: AWS service. Under \textbf{Use case},
  find and select \textbf{Personalize}
\item
  Attach the \texttt{AmazonS3ReadOnlyAccess} policy
\item
  Name it \texttt{PersonalizeS3Role} and choose \textbf{Create role}
\end{enumerate}

Add a bucket policy allowing Personalize to read. Create
\texttt{bucket-policy.json}:

\begin{Shaded}
\begin{Highlighting}[]
\FunctionTok{\{}
  \DataTypeTok{"Version"}\FunctionTok{:} \StringTok{"2012{-}10{-}17"}\FunctionTok{,}
  \DataTypeTok{"Statement"}\FunctionTok{:} \OtherTok{[}\FunctionTok{\{}
    \DataTypeTok{"Effect"}\FunctionTok{:} \StringTok{"Allow"}\FunctionTok{,}
    \DataTypeTok{"Principal"}\FunctionTok{:} \FunctionTok{\{} \DataTypeTok{"Service"}\FunctionTok{:} \StringTok{"personalize.amazonaws.com"} \FunctionTok{\},}
    \DataTypeTok{"Action"}\FunctionTok{:} \OtherTok{[}\StringTok{"s3:GetObject"}\OtherTok{,} \StringTok{"s3:ListBucket"}\OtherTok{]}\FunctionTok{,}
    \DataTypeTok{"Resource"}\FunctionTok{:} \OtherTok{[}
      \StringTok{"arn:aws:s3:::fcj{-}recsys{-}data{-}\textless{}your{-}name\textgreater{}"}\OtherTok{,}
      \StringTok{"arn:aws:s3:::fcj{-}recsys{-}data{-}\textless{}your{-}name\textgreater{}/*"}
    \OtherTok{]}
  \FunctionTok{\}}\OtherTok{]}
\FunctionTok{\}}
\end{Highlighting}
\end{Shaded}

\begin{Shaded}
\begin{Highlighting}[]
\ExtensionTok{aws}\NormalTok{ s3api put{-}bucket{-}policy }\AttributeTok{{-}{-}bucket}\NormalTok{ fcj{-}recsys{-}data{-}}\OperatorTok{\textless{}}\NormalTok{your{-}name}\OperatorTok{\textgreater{}}\NormalTok{ {-}{-}policy file://bucket{-}policy.json}
\end{Highlighting}
\end{Shaded}

\vspace{0.5em}\noindent\textbf{Step 3. Create a dataset
group}\par

\begin{enumerate}
\def\labelenumi{\arabic{enumi}.}
\tightlist
\item
  Open the
  \href{https://ap-southeast-1.console.aws.amazon.com/personalize/home?region=ap-southeast-1}{Amazon
  Personalize console}
\item
  Choose \textbf{Create dataset group}
\item
  \textbf{Name}: \texttt{fcj-recsys-dsg}, \textbf{Domain}: choose
  \textbf{Custom}
\item
  Choose \textbf{Create dataset group and continue}
\end{enumerate}

\begin{figure}
\centering
\pandocbounded{\includegraphics[keepaspectratio,alt={Create dataset group}]{5-Workshop/5.6/create-dsg.png}}
\caption{Create dataset group}
\end{figure}

\vspace{0.5em}\noindent\textbf{Step 4. Import the interactions
data}\par

\begin{enumerate}
\def\labelenumi{\arabic{enumi}.}
\tightlist
\item
  At the \textbf{Create interactions dataset} step, choose
  \textbf{Create dataset}
\item
  \textbf{Dataset name}: \texttt{interactions}
\item
  Choose \textbf{Create new schema}, name it
  \texttt{interactions-schema}, and paste:
\end{enumerate}

\begin{Shaded}
\begin{Highlighting}[]
\FunctionTok{\{}
  \DataTypeTok{"type"}\FunctionTok{:} \StringTok{"record"}\FunctionTok{,}
  \DataTypeTok{"name"}\FunctionTok{:} \StringTok{"Interactions"}\FunctionTok{,}
  \DataTypeTok{"namespace"}\FunctionTok{:} \StringTok{"com.amazonaws.personalize.schema"}\FunctionTok{,}
  \DataTypeTok{"fields"}\FunctionTok{:} \OtherTok{[}
    \FunctionTok{\{} \DataTypeTok{"name"}\FunctionTok{:} \StringTok{"USER\_ID"}\FunctionTok{,} \DataTypeTok{"type"}\FunctionTok{:} \StringTok{"string"} \FunctionTok{\}}\OtherTok{,}
    \FunctionTok{\{} \DataTypeTok{"name"}\FunctionTok{:} \StringTok{"ITEM\_ID"}\FunctionTok{,} \DataTypeTok{"type"}\FunctionTok{:} \StringTok{"string"} \FunctionTok{\}}\OtherTok{,}
    \FunctionTok{\{} \DataTypeTok{"name"}\FunctionTok{:} \StringTok{"EVENT\_TYPE"}\FunctionTok{,} \DataTypeTok{"type"}\FunctionTok{:} \StringTok{"string"} \FunctionTok{\}}\OtherTok{,}
    \FunctionTok{\{} \DataTypeTok{"name"}\FunctionTok{:} \StringTok{"TIMESTAMP"}\FunctionTok{,} \DataTypeTok{"type"}\FunctionTok{:} \StringTok{"long"} \FunctionTok{\}}
  \OtherTok{]}\FunctionTok{,}
  \DataTypeTok{"version"}\FunctionTok{:} \StringTok{"1.0"}
\FunctionTok{\}}
\end{Highlighting}
\end{Shaded}

\begin{enumerate}
\def\labelenumi{\arabic{enumi}.}
\setcounter{enumi}{3}
\tightlist
\item
  Choose \textbf{Next}, then \textbf{Import data directly into datasets}
\item
  Fill in:

  \begin{itemize}
  \tightlist
  \item
    \textbf{Import job name}: \texttt{import-interactions}
  \item
    \textbf{IAM role}: select \texttt{PersonalizeS3Role}
  \item
    \textbf{Data location}:
    \texttt{s3://fcj-recsys-data-\textless{}your-name\textgreater{}/personalize/interactions.csv}
  \end{itemize}
\item
  Choose \textbf{Start import}
\end{enumerate}

The import takes roughly 10 to 20 minutes. The status must reach
\textbf{Active}.

While waiting, take screenshots of the steps you have completed for your
report.

\vspace{0.5em}\noindent\textbf{Step 5. Train a solution}\par

\begin{enumerate}
\def\labelenumi{\arabic{enumi}.}
\tightlist
\item
  In the left menu, choose \textbf{Custom resources} then
  \textbf{Solutions and recipes}
\item
  Choose \textbf{Create solution}
\item
  Fill in:

  \begin{itemize}
  \tightlist
  \item
    \textbf{Solution name}: \texttt{fcj-user-personalization}
  \item
    \textbf{Solution type}: \textbf{Item recommendation}
  \item
    \textbf{Recipe}: \texttt{aws-user-personalization-v2}
  \end{itemize}
\item
  Choose \textbf{Create and train solution}
\end{enumerate}

Training takes roughly 40 to 90 minutes depending on data volume.

\vspace{0.5em}\noindent\textbf{Step 6. Review the evaluation
metrics}\par

Once the solution version becomes \textbf{Active}, open it to see the
metrics:

\begin{itemize}
\tightlist
\item
  \textbf{Precision@K}: share of correct items among the first K
  recommendations
\item
  \textbf{NDCG@K}: like Precision but position-aware; a correct item at
  the top counts more than one at the bottom
\item
  \textbf{MRR@K}: average position of the first correct recommendation
\item
  \textbf{Coverage}: share of the catalogue that ever gets recommended
\end{itemize}

\begin{figure}
\centering
\pandocbounded{\includegraphics[keepaspectratio,alt={Model metrics}]{5-Workshop/5.6/metrics.png}}
\caption{Model metrics}
\end{figure}

\vspace{0.5em}\noindent\textbf{Step 7. Create a campaign}\par

A campaign is the endpoint that serves real-time inference.

\begin{enumerate}
\def\labelenumi{\arabic{enumi}.}
\tightlist
\item
  Choose \textbf{Campaigns} then \textbf{Create campaign}
\item
  Fill in:

  \begin{itemize}
  \tightlist
  \item
    \textbf{Campaign name}: \texttt{fcj-recsys-campaign}
  \item
    \textbf{Solution}: pick the solution you just trained
  \item
    \textbf{Minimum provisioned TPS}: \texttt{1}
  \end{itemize}
\item
  Choose \textbf{Create campaign}
\end{enumerate}

A campaign starts \textbf{billing per hour} the moment it is created,
even with zero queries. It is the most expensive resource in this
workshop. Remember to delete it during clean-up.

\vspace{0.5em}\noindent\textbf{Step 8. Wire the campaign into
Lambda}\par

Copy the Campaign ARN, then add an environment variable to Lambda:

\begin{enumerate}
\def\labelenumi{\arabic{enumi}.}
\tightlist
\item
  Open the \texttt{fcj-recsys-api} function, \textbf{Configuration} tab,
  \textbf{Environment variables}
\item
  Choose \textbf{Edit} and add:

  \begin{itemize}
  \tightlist
  \item
    \textbf{Key}: \texttt{CAMPAIGN\_ARN}
  \item
    \textbf{Value}: the ARN you copied
  \end{itemize}
\item
  Choose \textbf{Save}
\end{enumerate}

\vspace{0.5em}\noindent\textbf{Verify}\par

\begin{Shaded}
\begin{Highlighting}[]
\ExtensionTok{aws}\NormalTok{ personalize{-}runtime get{-}recommendations }\DataTypeTok{\textbackslash{}}
  \AttributeTok{{-}{-}campaign{-}arn} \OperatorTok{\textless{}}\NormalTok{CAMPAIGN{-}ARN}\OperatorTok{\textgreater{}} \DataTypeTok{\textbackslash{}}
  \AttributeTok{{-}{-}user{-}id}\NormalTok{ user{-}001 }\DataTypeTok{\textbackslash{}}
  \AttributeTok{{-}{-}num{-}results}\NormalTok{ 5}
\end{Highlighting}
\end{Shaded}

A list of \texttt{itemId} values means the model is working. Try a few
different \texttt{user-id} values --- the results should differ, which
is your proof that recommendations are genuinely personalised.
